\documentclass[conference]{IEEEtran}

\usepackage{graphicx}
\usepackage{amsmath,amssymb}
\usepackage{float}
\usepackage{booktabs}
\usepackage{multirow}
\usepackage{xcolor}

\begin{document}

\title{
    \textbf{Benchmark Report: 502.gcc\_r} \\
    \Large\textbf{\textcolor{blue}{[ TRAIN INPUT ]}} \\
    \normalsize Register Spill Analysis on RISC-V via gem5
}

\author{
    \IEEEauthorblockN{Lütfullah Burak Kaya}
    \IEEEauthorblockA{
        Ozyegin University \\
        burak.kaya.50711@ozu.edu.tr
    }
    \and
    \IEEEauthorblockN{Ismail Akturk}
    \IEEEauthorblockA{
        Ozyegin University \\
        ismail.akturk@ozyegin.edu.tr
    }
}

\newcommand{\LongDate}{%
    \number\day\space
    \ifcase\month
    \or January\or February\or March\or April\or May\or June%
    \or July\or August\or September\or October\or November\or December%
    \fi
    \space \number\year
}

\twocolumn[
\begin{@twocolumnfalse}
\maketitle
\begin{center}{\small \LongDate}\end{center}
\vspace{1em}
\end{@twocolumnfalse}
]

% ================================================

\begin{abstract}
This report presents the register spill characterization of the
\texttt{502.gcc\_r} benchmark from SPEC CPU2017, executed on a
RISC-V (RV64GC) target using the gem5 TimingSimpleCPU simulator
with the \textbf{train} input dataset. A custom SpillDetector was
integrated into gem5 to count store-then-load patterns on the stack
within a user-defined Region of Interest (ROI). The train dataset
consists of three C source files (\texttt{200.c}, \texttt{scilab.c},
\texttt{train01.c}), each compiled with \texttt{-O3 -finline-limit=50000}
in a separate gem5 invocation. The combined spill rate across all
three compilations is \textbf{8.12\%}, the highest of all 17 evaluated
SPEC CPU2017 benchmarks, corresponding to 17.9 billion spill events
across 220.8 billion ROI instructions. Two-thirds of all store
instructions (66.7\%) are spill writes, and one-third of all load
instructions (34.2\%) are spill reloads. Compared to the test input
(6.92\%), the train input reveals a higher and more stable spill rate,
confirming that \texttt{gcc\_r} is the most register-pressure
intensive workload in the SPEC CPU2017 suite on RISC-V RV64GC.
\end{abstract}

% =========================================================
\section{Benchmark Overview}

\texttt{502.gcc\_r} is the SPEC CPU2017 rate variant of GCC~4.3.2,
the GNU C compiler. The benchmark executes a RISC-V binary of GCC
that produces \textbf{x86} assembly output: the compiler itself runs
on RISC-V inside gem5, but it targets the x86 ISA. This is expected
behavior --- the output \texttt{.s} files contain x86 assembly, while
all of the register pressure measured by the SpillDetector arises from
the RISC-V execution of the compiler itself.

GCC's internal architecture involves a large number of passes operating
on tree and RTL intermediate representations, each traversing and
transforming complex recursive data structures. Function calls between
passes are numerous and pass deep trees of live IR pointers, creating
sustained register pressure throughout the compilation pipeline.

The RISC-V binary (\texttt{gcc\_r.riscv}) was compiled with gem5
ROI markers inserted around the \texttt{toplev\_main(argc, argv)}
call in \texttt{src/main.c}:
\begin{itemize}
    \item \texttt{m5\_work\_begin(0,0)} --- before \texttt{toplev\_main()}
    \item \texttt{m5\_work\_end(0,0)}   --- after  \texttt{toplev\_main()}
\end{itemize}
This ROI encompasses the complete compilation of each input C file.

\subsection{Input Datasets}

SPEC CPU2017 provides three dataset sizes for \texttt{gcc\_r}
(Table~\ref{tab:inputs}). The train dataset consists of three
separate C files, each compiled in a separate invocation.

\begin{table}[H]
\centering
\caption{gcc\_r Input Dataset Sizes}
\label{tab:inputs}
\begin{tabular}{lll}
\toprule
\textbf{Dataset} & \textbf{Input files} & \textbf{Total size} \\
\midrule
test  & \texttt{t1.c} (hello world)  & 69\,B \\
train & \texttt{200.c}, \texttt{scilab.c}, \texttt{train01.c} & $\approx$4\,MB \\
ref   & \texttt{gcc-pp.c}, \texttt{gcc-smaller.c}, \texttt{ref32.c} & $\approx$6\,MB \\
\bottomrule
\end{tabular}
\end{table}

The \textbf{train} dataset was used for this study. Each of the three
C files is compiled independently with identical flags
(\texttt{-O3 -finline-limit=50000}), as specified by the SPEC
\texttt{control} file. Each compilation was run as a separate gem5
simulation with its own output directory:

\begin{itemize}
    \item \texttt{200.c} (1.81\,MB, 63,888 lines) $\rightarrow$ \texttt{m5out\_train\_200}
    \item \texttt{scilab.c} (1.09\,MB, 58,762 lines) $\rightarrow$ \texttt{m5out\_train\_scilab}
    \item \texttt{train01.c} (1.19\,MB, 50,779 lines) $\rightarrow$ \texttt{m5out\_train\_train01}
\end{itemize}

% =========================================================
\section{Simulation Setup}

\subsection{gem5 Configuration}

\begin{itemize}
    \item \textbf{CPU model:}   TimingSimpleCPU
    \item \textbf{ISA:}         RISC-V RV64GC
    \item \textbf{Caches:}      enabled (default L1 I/D + L2)
    \item \textbf{Memory:}      4\,GB simulated physical RAM
    \item \textbf{Spill mode:}  summary only (\texttt{spill\_verbose=False})
\end{itemize}

Three gem5 invocations were launched in parallel, one per input file:

\begin{verbatim}
./build/RISCV/gem5.opt
  --outdir=benchmarks/gcc/m5out_train_200
  configs/deprecated/example/se.py
  --cmd=benchmarks/gcc/gcc_r.riscv
  --options="benchmarks/gcc/200.c
    -O3 -finline-limit=50000
    -o benchmarks/gcc/200.s"
  --cpu-type=TimingSimpleCPU --caches --mem-size=4GB
\end{verbatim}

The same invocation pattern was used for \texttt{scilab.c} and
\texttt{train01.c} with their respective output directories.

% =========================================================
\section{Simulation Results}

All three compilations completed successfully, each producing an
x86 assembly output file. The results are presented per input file
and as a combined aggregate.

\subsection{Pre-ROI Context}

All three runs share the same binary and startup path, producing
identical pre-ROI counters at ROI entry (Table~\ref{tab:preroi}).

\begin{table}[H]
\centering
\caption{Global Counters at ROI Entry (identical for all 3 runs)}
\label{tab:preroi}
\begin{tabular}{lr}
\toprule
\textbf{Metric} & \textbf{Value} \\
\midrule
Total instructions (pre-ROI) & 2,391,942 \\
Total spills detected (pre-ROI) & 79,848 \\
Pre-ROI spill rate & 3.34\% \\
\bottomrule
\end{tabular}
\end{table}

\subsection{Per-File ROI Spill Statistics}

Table~\ref{tab:perfile} shows the spill statistics for each
individual compilation.

\begin{table}[H]
\centering
\caption{Per-File ROI Spill Statistics}
\label{tab:perfile}
\begin{tabular}{lrrr}
\toprule
\textbf{Input} & \textbf{ROI Instr} & \textbf{ROI Spills} & \textbf{Rate} \\
\midrule
200.c     & 118.4\,B & 9,876\,M & 8.3381\% \\
scilab.c  &  92.3\,B & 7,232\,M & 7.8336\% \\
train01.c &  10.0\,B &   811\,M & 8.0836\% \\
\midrule
\textbf{Combined} & \textbf{220.8\,B} & \textbf{17,918\,M} & \textbf{8.1160\%} \\
\bottomrule
\end{tabular}
\end{table}

\subsection{Combined ROI Spill Statistics}

Table~\ref{tab:roi} presents the aggregate statistics across all
three compilations.

\begin{table}[H]
\centering
\caption{Combined ROI Spill Statistics --- 502.gcc\_r (train)}
\label{tab:roi}
\begin{tabular}{lr}
\toprule
\textbf{Metric} & \textbf{Value} \\
\midrule
ROI instructions (total)   & 220,789,486,676 \\
ROI stores (total)         &  26,867,426,671 \\
ROI loads (total)          &  52,398,708,738 \\
\midrule
\textbf{ROI spills (total)} & \textbf{17,918,365,207} \\
\textbf{Combined spill rate} & \textbf{8.1160\%} \\
\bottomrule
\end{tabular}
\end{table}

\subsection{Timing}

\begin{table}[H]
\centering
\caption{Per-File Simulation Timing}
\label{tab:timing}
\begin{tabular}{lrr}
\toprule
\textbf{Input} & \textbf{Sim time} & \textbf{Host time} \\
\midrule
200.c     & 293.45\,s & 77,003\,s ($\approx$21.4\,h) \\
scilab.c  & 232.52\,s & 60,849\,s ($\approx$16.9\,h) \\
train01.c &  26.47\,s &  6,475\,s ($\approx$1.8\,h) \\
\midrule
\textbf{Total} & \textbf{552.43\,s} & \textbf{144,327\,s ($\approx$40.1\,h)} \\
\bottomrule
\end{tabular}
\end{table}

% =========================================================
\section{ROI Load and Store Breakdown}

\subsection{Load and Store Composition}

Within the 220.8 billion ROI instructions, 26.9 billion are stores
and 52.4 billion are loads. Table~\ref{tab:lsbreakdown} decomposes
each category into spill-related and non-spill traffic.

\begin{table}[H]
\centering
\caption{ROI Load/Store Breakdown --- 502.gcc\_r (train)}
\label{tab:lsbreakdown}
\begin{tabular}{lrr}
\toprule
\textbf{Category} & \textbf{Count} & \textbf{Share} \\
\midrule
\multicolumn{3}{l}{\textit{Store instructions}} \\
\quad Spill stores (writes to stack)    & 17,918,365,207 & 66.69\% of stores \\
\quad Non-spill stores                  &  8,949,061,464 & 33.31\% of stores \\
\quad \textbf{Total stores}             & \textbf{26,867,426,671} & 12.17\% of ROI instr \\
\midrule
\multicolumn{3}{l}{\textit{Load instructions}} \\
\quad Spill reloads (reads from stack)  & 17,918,365,207 & 34.19\% of loads \\
\quad Non-spill loads                   & 34,480,343,531 & 65.81\% of loads \\
\quad \textbf{Total loads}              & \textbf{52,398,708,738} & 23.73\% of ROI instr \\
\bottomrule
\end{tabular}
\end{table}

The key observation is that \textbf{two-thirds of all store operations
are spill writes}. This means the stack dominates write traffic: for
every 3 store instructions executed by GCC on RISC-V, 2 are spilling
a register to memory. On the load side, 1 in 3 loads is a spill
reload. The asymmetry (66.7\% vs.\ 34.2\%) arises because loads
outnumber stores by $1.95\times$ in GCC's IR-traversal-dominated
execution, while each spill event contributes exactly one store and
one reload.

\subsection{Spill Read/Write Decomposition}

Every register spill is a two-part operation:
\begin{enumerate}
    \item \textbf{Spill write (store):} the register value is written
          to the stack frame to free the physical register.
    \item \textbf{Spill reload (load):} the value is read back from
          the stack when the register is needed again.
\end{enumerate}

The SpillDetector counts a spill event when it observes a
store-then-load pair targeting the same stack address. Each counted
spill therefore contributes exactly one store and one load to the
memory traffic. Table~\ref{tab:rwdecomp} summarises the split.

\begin{table}[H]
\centering
\caption{Spill Read/Write Decomposition --- 502.gcc\_r (train)}
\label{tab:rwdecomp}
\begin{tabular}{lrr}
\toprule
\textbf{Spill component} & \textbf{Count} & \textbf{Fraction of total} \\
\midrule
Spill writes (stores) & 17,918,365,207 & 66.69\% of ROI stores \\
Spill reloads (loads) & 17,918,365,207 & 34.19\% of ROI loads  \\
\midrule
Combined spill memory ops & 35,836,730,414 & --- \\
Total ROI memory ops & 79,266,135,409 & --- \\
\textbf{Spill memory fraction} & --- & \textbf{45.22\%} \\
\bottomrule
\end{tabular}
\end{table}

Nearly half (45.2\%) of all memory operations in the ROI are
spill-related. The remaining 54.8\% are genuine data loads and
stores --- pointer dereferences into GCC's IR tree, output writes
to the assembly buffer, and internal bookkeeping. The high spill
memory fraction is a direct consequence of RISC-V's 32-register
file being exhausted by GCC's deep, pointer-rich pass call chains.

% =========================================================
\section{Analysis}

\subsection{Spill Rate Interpretation}

The combined spill rate of \textbf{8.12\%} places \texttt{gcc\_r}
(train) first among all 17 evaluated SPEC CPU2017 benchmarks
(Table~\ref{tab:compare}). The table includes three additional
columns: spill rate as a fraction of ROI loads (Spill/Load) and
as a fraction of ROI stores (Spill/Store), revealing how spill
pressure manifests differently in the load and store streams.

\begin{table*}[t]
\centering
\caption{Cross-Benchmark Spill Rate Comparison (Test Input Unless Noted)}
\label{tab:compare}
\begin{tabular}{lrrrr}
\toprule
\textbf{Benchmark} & \textbf{ROI Spills} & \textbf{Spill Rate}
    & \textbf{Spill/Load} & \textbf{Spill/Store} \\
\midrule
502.gcc\_r \textit{(train)}  & 17,918,365,207 & \textbf{8.1160\%} & 34.19\% & 66.69\% \\
500.perlbench\_r             &     12,728,325,675 & 8.5557\% & 31.24\% & 54.75\% \\
507.cactuBSSN\_r             &  3,782,099,438 & 7.8756\% & 30.31\% & 54.66\% \\
511.povray\_r                &    174,359,922 & 7.7809\% & 24.34\% & 60.54\% \\
520.omnetpp\_r               &    894,619,681 & 7.1994\% & 32.02\% & 59.98\% \\
502.gcc\_r \textit{(test)}   &      1,108,838 & 6.9185\% & 37.49\% & 56.70\% \\
523.xalancbmk\_r             &     18,381,250 & 4.6427\% & 17.34\% & 47.14\% \\
541.leela\_r                 &  1,158,236,271 & 4.4438\% & 22.66\% & 65.94\% \\
519.lbm\_r                   &    274,130,440 & 4.2130\% & 21.47\% & 35.34\% \\
531.deepsjeng\_r             &  1,928,903,656 & 4.1749\% & 24.14\% & 56.81\% \\
525.x264\_r                  &    767,981,981 & 3.2812\% & 15.87\% & 31.98\% \\
544.nab\_r                   &    151,310,616 & 2.2612\% &  9.93\% & 47.06\% \\
538.imagick\_r               &      2,424,899 & 2.0601\% & 11.77\% & 26.99\% \\
557.xz\_r                    &     38,411,428 & 1.9069\% & 10.44\% & 20.93\% \\
505.mcf\_r                   &    446,113,136 & 1.8263\% & --- & --- \\
508.namd\_r                  &    327,718,955 & 1.0269\% &  3.91\% & 16.84\% \\
510.parest\_r                &    127,688,196 & 0.3427\% &  1.07\% & 12.64\% \\
\bottomrule
\end{tabular}
\end{table*}

The Spill/Store column is consistently higher than Spill/Load across
all benchmarks because stores are fewer than loads in typical programs
(most computation is read-dominated). Benchmarks with a high
Spill/Store ratio --- such as \texttt{gcc\_r} train (66.7\%),
\texttt{leela\_r} (65.9\%), and \texttt{povray\_r} (60.5\%) ---
are dominated by call-heavy workloads where the register file is
repeatedly exhausted at function boundaries. Benchmarks with a
lower Spill/Store ratio --- such as \texttt{lbm\_r} (35.3\%) and
\texttt{x264\_r} (32.0\%) --- have more data-intensive store traffic
(array writes, output buffers) that dilutes the spill fraction of
the store stream.

\subsection{Train vs.\ Test Input Comparison}

The train dataset raises the spill rate from 6.92\% (test) to
8.12\% (train), a relative increase of 17.3\%. The test input
(\texttt{t1.c}, 69\,B) is a trivially small program that GCC
compiles in 16M ROI instructions, exercising only a narrow
subset of GCC's optimisation passes. The train input files are
four to five orders of magnitude larger (up to 1.81\,MB), causing
GCC to invoke the full set of GIMPLE and RTL optimisation passes
including inlining, loop transforms, instruction scheduling, and
register allocation. These passes traverse deeply nested IR trees,
generating the sustained register pressure that the test input
does not expose.

The Spill/Load ratio drops from 37.5\% (test) to 34.2\% (train),
while Spill/Store rises from 56.7\% (test) to 66.7\% (train).
The larger train files generate more IR node construction and
transformation (additional stores), of which a growing fraction
are spills, while the relative increase in genuine data loads
(pointer reads through larger IR trees) partially dilutes the
spill reload fraction.

\subsection{ROI Coverage}

Across all three runs, the ROI covers $220{,}789{,}486{,}676$ out
of $220{,}801{,}545{,}416$ total simulated instructions:
\[
\frac{220{,}789{,}486{,}676}{220{,}801{,}545{,}416} \approx 99.99\%
\]
The near-perfect ROI coverage confirms that \texttt{toplev\_main()}
encompasses virtually all benchmark computation in every invocation.

% =========================================================
\section{Conclusion}

The \texttt{502.gcc\_r} benchmark with the train input exhibits
a combined register spill rate of 8.12\% across three C source
file compilations on RISC-V RV64GC, the highest of all 17 evaluated
SPEC CPU2017 benchmarks. With 17.9 billion spill events across
220.8 billion ROI instructions and a 99.99\% ROI coverage, the
benchmark provides a comprehensive characterisation of GCC's
register behaviour at scale.

The load/store breakdown reveals that spills account for 66.7\%
of all store traffic and 34.2\% of all load traffic, collectively
consuming 45.2\% of all memory operations in the ROI. This confirms
that stack spill traffic is the dominant form of memory activity in
GCC's RISC-V execution --- not data structure reads or output writes,
but register save/restore sequences driven by call-chain depth and
the 32-register architectural limit. The 17.3\% spill-rate increase
from test to train input demonstrates that the train dataset exercises
the full GCC optimisation pipeline and yields a more representative
characterisation of real-world compilation workloads.

\end{document}
